%
% Core CFD solver
% -------------
\section{Core CFD solver}
\label{sec-corecfd}

The core CFD solver is written in C with a Python wrapper written using ctypes.
The C code is self contained and does not rely on any third party libraries other
than those provided by the compiler. e.g. libm on Linux.


%
% C Build
% -------
\subsection{Building the C code}
\label{subsec-buildc}


\subsubsection{Linux based systems}
The C code can be built as a stand-alone executable or as a library to be linked to a Python code.

\begin{itemize}
\item To make the C executable, type {\tt make} in the {\tt src} directory.
The executable {\tt nozzle} is placed in the {\tt bin} directory.

\item To make the C library, type {\tt make lib} in the {\tt src} directory.
The shared object library {\tt libnozzle.so} is placed in the {\tt lib} directory.

\item To rebuild from scratch , type  {\tt make clean} in the {\tt src} directory.
\end{itemize}


\subsubsection{Windows based systems}

Not cuurently supported.



%
% Running the code
% ----------------
\subsection{Running the code}
\label{subsec-run}

There are 3 modes of running the code include in the distribution.

\begin{itemize}
\item To run the C-code directly type: {\tt [path]/bin/nozzle -i [path]inputfile.xml}.

      Typing {\tt [path]/bin/nozzle -h} gives a short list of the input options.

      Sample input files are included in the {\tt input} directory of the distribution.
      Note each input file can be run in the four possible combinations of incompressible/compressible flow and
      SIMPLE/coupled solvers. Which combination is run is determined by the {\tt setup} section of the
      input file as described in \Cref{subsec-ifile}.

\item To run the C solver as a library call from Python type: {\tt [path]/script/run\_cprog.py -i [path]inputfile.xml}.

      This simply demonstrates the infrastructure used to call the C library from within a Python script.

\item To use the C library to set-up matrices to be solved using the {\tt scipy} linear algebra routine type:
      {\tt [path]/script/run\_pysolve.py -i [path]inputfile.xml}.  
\end{itemize}

The final option provides a template for replacing the {\tt scipy} linear solver with a quantum linear equation solver
as described in \Cref{subsec-qles}.


%
% Python C interface
% ------------------
\subsection{Python-C interface}
\label{subsec-py2c}
There are two layers to the Python-C interface:
\begin{itemize}
  \item Data instantiation when python is the main routine.
  \item Data exchange to {\it get} a matrix and a right-hand side vector, 
        and to {\it put} a solution vector.
\end{itemize}

Both of these depend on sharing C structures with the Python code.


%
% Data instantiation
% ------------------
\subsubsection{Data instantiation}
\label{subsubsec-data-inst}

\Cref{list-noz-cstruct} shows the high level structure used by the C solver.
As can be seen, this is a structure of structures which, collectively, contain all the
data used by the solver. These are not further defined here as they are not needed by
the Python layer.

\begin{lstlisting}[numbers=left, language=C, caption=C master structure for nozzle data, label=list-noz-cstruct]
typedef struct
{
        SETUP   *setup;
        MESH    *mesh;
        FLOW    *flow;
        SOLVER  *solver;
        FILES   *files;
        ARRAYS  *arrays;
        HISTORY *history;
} NOZZLE;
\end{lstlisting}

In fact, the Python layer needs only manage the pointer which it can do by treating it
as a {\tt void}. The only data passed between the C and Python layers is the memory address of
the nozzle pointer.
\Cref{list-noz-pyinst} gives code segments which call C routines to instantiate an empty nozzle structure;
load data from an input file; and, free the structure after the calculations have been completed.
Note that the call to {\tt c\_lib.c2py\_noz\_create} on line 2 passes the pointer by reference. 
This is needed because the memory for the pointer is allocated with the C code, so it must return
a pointer to a pointer.
The variable {\tt i\_cptr} on line 4 is the name of the input file that is read as a command 
line option.


\begin{lstlisting}[numbers=left, language=Python, caption=Python instantiation., label=list-noz-pyinst]
    noz_cptr = c_void_p()
    c_lib.c2py_noz_create(byref(noz_cptr))
    ...
    c_lib.c2py_read_input(i_cptr, noz_cptr)
    ...
    c_lib.c2py_noz_free(noz_cptr)
\end{lstlisting}

%
% Data exchange
% -------------
\subsubsection{Data exchange}
\label{subsubsec-data-exch}

In order to investigate quantum linear solvers, it is necessary to share matrices and vectors
between the C and Python layers. This is done by using a common structure definition.
\Cref{list-smat-cstruct} shows the C structure and \Cref{list-smat-pystruct} shows the
corresponding Python structure. Both used the same compressed row format.

\begin{lstlisting}[numbers=left, language=C, caption=C sparse matrix structure., label=list-smat-cstruct]
typedef struct
{
        int32_t nr;
        int32_t nc;
        int32_t nnz;
        int32_t *col;
        int32_t *rst;
        double  *val;
        double  *b;
} SMATRIX;
\end{lstlisting}

Note that integers are declared as {\tt int32} to ensure both the C and Python layers use 
32 bit integers.
Note that the sparse matrix structures also includes the residual, or right hand side vector, {\tt b}.
This is a convenience as both the matrix and vector can be retrieved in a single call.
It is also convenient on the C side as the matrix entries and residuals are computed at the
same time.
The solution vector is computed separately and, therefore, not included in the structure.


\begin{lstlisting}[numbers=left, language=Python, caption=Python sparse matrix structure., label=list-smat-pystruct]
class SMATRIX(Structure):
      _fields_ = [("nr",  c_int32),
                  ("nc",  c_int32),
                  ("nnz", c_int32),
                  ("col", POINTER(c_int32)),
                  ("rst", POINTER(c_int32)),
                  ("val", POINTER(c_double)),
                  ("b",   POINTER(c_double))
                 ]
      _defaults_ = {"nr" : 0,
                    "nc" : 0,
                    "nnz": 0,
                    "col": None,
                    "rst": None,
                    "val": None,
                    "b"  : None}
\end{lstlisting}


\Cref{list-get-smat} shows how to {\it get} a matrix and vector from the C code and
convert them to a scipy sparse matrix and numpy vector.
This is taken from the script {\tt run\_pysolve.py}
On line 10, the {\tt scipy.sparse.linalg} spsolve routine is used to solve for {\tt x}.
Lines 12-13 show how to {\it put} the solution vector back into the C code.


\begin{lstlisting}[numbers=left, language=Python, caption=Code segment for Python-C matrix and vector exchange., label=list-get-smat]
    S = SMATRIX()
    c_lib.c2py_get_smat(noz_cptr, byref(S))

    col = np.ctypeslib.as_array(S.col, shape=(S.nnz,))
    rst = np.ctypeslib.as_array(S.rst, shape=(S.nr+1,))
    val = np.ctypeslib.as_array(S.val, shape=(S.nnz,))
    b   = np.ctypeslib.as_array(S.b,   shape=(S.nr,))

    SP = sparse.csr_matrix((val, col, rst),shape=(S.nr, S.nc), dtype=np.double)
    x  = spsolve(SP, b)

    xc = x.ctypes.data_as(POINTER(c_double))
    rc = c_lib.c2py_put_soln(noz_cptr, xc)
\end{lstlisting}

As can be seen, the above code segments are independent of the CFD flow type and solution scheme.
These are contained within the C layer and controlled via a single input file.



%
% QLES solver
% ----------
\subsection{Linking to a quantum linear equation solver}
\label{subsec-qles}

In principle, linking to a quantum solver is {\it just} a matter of replacing the call to
{\tt spsolve} in line 10 of \Cref{list-get-smat} with a call to the equivalent quantum solver.
However, the quantum solver is likely to need a lot of additional input, including additional
data such as phase factors for QSVT.
The user has a choice to hardwire them for initial development testing, add them as additonal 
input options to the python script, or include them in an auxillary input file.

Note that the call to {\tt spsolve} is within an iterative loop. Only data that changes as part of the
non-linear outer iterations should be included within the loop.
Data input and set-up should be outside the loop.

It is recommended to make a copy of the file {\tt run\_pysolve.py} and use that as a template 
for the QLES solver.


%
% Input file
% ----------
\subsection{Input file}
\label{subsec-ifile}

There are sample input files in the {\tt input} directory of the distribution.
The files are structured to include options for all combinations of flow type and solver.
The choice to be run is controlled by the {\tt setup} section as described in \Cref{subsubsec-ifile-overview}.
If files need to be simplified, the sections for the unused options can be removed. 

%
% overview
% --------
\subsubsection{Overview}
\label{subsubsec-ifile-overview}

The input file has been designed for all 4 combinations of flow regime and
solver type to be held in a single file with two binary flags to select which option is executed.
This adds some verbosity to the file but eases the number of files that need to be managed.
This does not prevent separate files being used for separate cases if that is preferred.
The high-level structure of the file is shown in  \Cref{list-xml01}.

\begin{lstlisting}[numbers=left, language=xml, caption=Input file main sections., label=list-xml01]
<?xml version="1.0" encoding="UTF-8"?>
<nozzle>
  <setup>
    <casename>nozzle</casename>
    <compressible>0</compressible>
    <couple>1</couple>
  </setup>

  <mesh>...</mesh>

  <incompressible>
    <flow>...</flow>
    <solver>
      <simple>...</simple>
      <couple>...</couple>
    </solver>
  </incompressible>

  <compressible>
    <flow>...</flow>
    <solver>
      <simple>...</simple>
      <couple>...</couple>
    </solver>
  </compressible>
</nozzle>
\end{lstlisting}


The {\tt setup} section contains a case name and binary variables to
toggle the flow and solver types. 
This section in \Cref{list-xml01} specifies incompressible flow using the coupled solver.
See \Cref{subsubsec-naming} for how the case name is used to set the names of the 
output files based on the setup.

The {\tt mesh} section is common to all setup combinations and is described below.
There are separate sections for incompressible and compressible flow with
each containing {\tt flow} and {\tt solver} sections.
The solver sections contain separate {\tt simple} and {\tt couple} sections for
each solver type. As shown below, the variables in each section depend on the
set-up and where variables have the same name they generally have different values,
according to the section they are in. 

All variables have default values. Those variables that use the default values can
be omitted from the input file for brevity, if desired.


%
% mesh
% ----
\subsubsection{Mesh data}

\Cref{list-xml02} shows the {\tt mesh} section of the input file.

\begin{lstlisting}[numbers=left, language=xml, caption=Input file mesh section., label=list-xml02]
  <mesh>
    <nx>16</nx>
    <ain>1.0</ain>
    <da>-0.1</da>
  </mesh>
\end{lstlisting}

The variables are:
\begin{itemize}
  \item{\tt nx} -  the number of axial stations in the mesh, this value is used as part
                   of the output file names.
  \item{\tt ain} - for incompressible flows this is the inlet area of the channel.
                   For compressible flows this is a scale factor for the inlet area,
                   which  can be used to scale the matrix entries but is usually set to {\tt 1.0}.
  \item{\tt da}  - the area contraction at the midpoint of the channel for incompressible flows.
                   This is specified as a fraction of the inlet area and is negative for an
                   area contraction and positive for an area expansion.
\end{itemize}


%
% incomp-simple
% -------------
\subsubsection{Incompressible flow, SIMPLE solver}
\label{subsubsec-insi-xml}

\Cref{list-xml03} shows the section of the input file for incompressible flow with the SIMPLE solver.
Note that all flow variables are non-dimensional.

\begin{lstlisting}[numbers=left, language=xml, caption=Input file for incompressible + SIMPLE solver., label=list-xml03]
  <incompressible>
    <flow>
      <rho>1.0</rho>                   <!-- constant density -->
      <bc>
        <ptin>1.5</ptin>               <!-- inlet total pressure -->
        <pex>0.95</pex>                <!-- exit static pressure -->
      </bc>
    </flow>

    <solver>
      <simple>
        <niters>2000</niters>          <!-- number of iterations -->
        <relaxu>1.0</relaxu>           <!-- u-relaxation factor -->
        <relaxp>1.0</relaxp>           <!-- p-relaxation factor -->
        <toleru>1.0e-12</toleru>       <!-- u-convergence tolerance -->
        <tolerp>1.0e-12</tolerp>       <!-- p-convergence tolerance -->
      </simple>
    </solver>
  </incompressible>
\end{lstlisting}

The variables are defined in \Cref{list-xml03} and should be self-explanatory.
As shown in \Cref{tab-pc-incomp}, the number of iterations varies considerably
for different mesh sizes. 
{\tt niters} is the maximum number of iterations, the calculation will terminate
when it reaches either {\tt niters} or both of the convergence tolerances have been meet.


%
% incomp-couple
% -------------
\subsubsection{Incompressible flow, coupled solver}
\label{subsubsec-inco-xml}

\Cref{list-xml04} shows the section of the input file for incompressible flow with the coupled solver.
Note that all flow variables are non-dimensional.

\begin{lstlisting}[numbers=left, language=xml, caption=Input file for incompressible + coupled solver., label=list-xml04]
  <incompressible>
    <flow>...  </flow>

    <solver>
      <couple>
        <niters>10</niters>            <!-- number of iterations -->
        <presim>0</presim>             <!-- number SIMPLE iterations before coupled solver -->
        <relaxc>1.0</relaxc>           <!-- relaxation factor applied to all variables -->
        <rrampc>0</rrampc>             <!-- number of iterations over which to ramp relaxc -->
        <tolerc>1.0e-12</tolerc>       <!-- convergence tolerance for all variables -->
      </couple>
    </solver>
  </incompressible>
\end{lstlisting}

The {\tt flow} is not included as it is identical to \Cref{list-xml03}.
The number of iterations is as before, but has a much smaller value as \Cref{tab-imp-comp}
shows that number of iterations is less than 10 for all mesh sizes.
A single convergence tolerance, {\tt tolerc}, is used for the coupled set of equations.
The variables {\tt presim}, {\tt relaxc} and {\tt rrampc} are needed for the compressible 
coupled solver and are described in \Cref{subsubsec-coco-xml}. 
The values of these variables in \Cref{list-xml04} are the defaults and these can be omitted
from the input file for incompressible flows.


%
% comp-simple
% -------------
\subsubsection{Compressible flow, SIMPLE solver}
\label{subsubsec-cosi-xml}

\Cref{list-xml05} shows the section of the input file for compressible flow with the SIMPLE solver.
Note that all flow variables are non-dimensional.

\begin{lstlisting}[numbers=left, language=xml, caption=Input file for incompressible + SIMPLE solver., label=list-xml05]
  <compressible>
    <flow>
      <bc>
        <ptin>1.5</ptin>               <!-- inlet total pressure -->
        <min>0.8</min>                 <!-- inlet Mach number -->
        <pex>0.75</pex>                <!-- exit static pressure -->
      </bc>
      <gaslaw>
        <gamma>1.4</gamma>             <!-- ratio of specific heats -->
        <cp>3.5</cp>                   <!-- specific heat at constant pressure -->
        <rgas>1.0</rgas>               <!-- perfect gas constant -->
        <upwind>2</upwind>             <!-- 1, 2,3 for 1st, 2nd and 3-pt density upwinding -->
      </gaslaw>
    </flow>

    <solver>
      <simple>
        <niters>5000</niters>          <!-- number of iterations -->
        <relaxu>1.0</relaxu>           <!-- u-relaxation factor -->
        <relaxp>1.0</relaxp>           <!-- p-relaxation factor -->
        <toleru>1.0e-12</toleru>       <!-- u-convergence tolerance -->
        <tolerp>1.0e-12</tolerp>       <!-- p-convergence tolerance -->
      </simple>
    </solver>
  </compressible>
\end{lstlisting}

%
% comp-couple
% -------------
\subsubsection{Compressible flow, coupled solver}
\label{subsubsec-coco-xml}

\Cref{list-xml06} shows the section of the input file for compressible flow with the coupled solver.
Note that all flow variables are non-dimensional.

\begin{lstlisting}[numbers=left, language=xml, caption=Input file for incompressible + coupled solver., label=list-xml06]
  <compressible>
    <flow>...</flow>

    <solver>
      <couple>
        <niters>2000</niters>          <!-- number of iterations -->
        <presim>10</presim>            <!-- number SIMPLE iterations before coupled solver -->
        <relaxc>0.75</relaxc>          <!-- relaxation factor applied to all variables -->
        <rrampc>20</rrampc>            <!-- number of iterations over which to ramp relaxc -->
        <tolerc>1.0e-12</tolerc>       <!-- convergence tolerance for all variables -->
      </couple>
    </solver>
  </compressible>
\end{lstlisting}





%
% Output file
% -----------
\subsection{Output files}
\label{subsec-ofile}

\subsubsection{File naming convention}
\label{subsubsec-naming}

There are two output files: a convergence history and a flow solution file. 
The naming convention for both files is \textit{casename\_nx\_ffss.ext}, where:
\begin{itemize}
\item \textit{casename} is the user specified name in the input file.
\item \textit{nx} is the number of stations.
\item \textit{ff} is a two character identifier for the flow type: \textit{in} for
      incompressible flow or \textit{co} for compressible flow.
\item \textit{ss} is a two character identifier for the solver type: \textit{si} for
      the SIMPLE solver or \textit{co} for the coupled solver.
\item \textit{ext} is the three character file extension: \textit{hst} for the convergence
      history file and \textit{sol} for the solution file.
\end{itemize}

For example, the history file for an incompressible SIMPLE case named 'nozzle with 16 stations is
{\tt nozzle\_016\_insi.hst}.
If cases with more that 999 stations are needed, modify the routine {\tt noz\_files} to change the
format statement in the {\tt snprintf} statements to increase '3' to the number of digits needed.



\subsubsection{Convergence history}
\label{subsubsec-form-conv}

The contents of the convergence history file depend on the flow type.
In both cases the file contains the root mean square (RMS) of the changes in the flow variables between
the current iteration and the previous one; and, the residual errors, e.g $R^u$, of the flow equations
at the start of the iteration.
Both should reduce towards the convergence tolerance as the iterations proceed.
For incompressible cases, the history file contains the RMS changes for $u$ and $p$ and the RMS resisuals
of the momentum and continuity equations.
For compressible flows, the RMS changes in density and residual of the gas equation are also included.
Sample files for both flow types are listed below, showing the first 9 iterations.

See \Cref{subsec-res-conv} for details on how to plot the convergence history of a single run and
to compare the convergence histories of multiple runs.

\begin{verbatim}
head nozzle_016_insi.hst
"iter",  "du_rms",        "dp_rms",        "resu_rms",      "resp_rms"
 1,  0.00000000e+00,  9.95624279e-02,  0.00000000e+00,  1.62406810e-02
 2,  3.42680753e-02,  9.21423809e-02,  3.56263102e-02,  2.32527299e-02
 3,  4.09420463e-03,  1.56536726e-03,  4.31829249e-03,  1.56084554e-03
 4,  3.82501330e-03,  2.49477798e-03,  4.06783477e-03,  8.77342756e-04
 5,  3.60318337e-03,  2.36292698e-03,  3.84504091e-03,  8.24039930e-04
 6,  3.39492043e-03,  2.23327519e-03,  3.63437959e-03,  7.76310997e-04
 7,  3.19933160e-03,  2.11074512e-03,  3.43527772e-03,  7.31505443e-04
 8,  3.01557136e-03,  1.99495898e-03,  3.24709921e-03,  6.89419128e-04
 9,  2.84285931e-03,  1.88554307e-03,  3.06924290e-03,  6.49871476e-04


head nozzle_016_coco.hst
"iter", "du_rms",    "dp_rms",     "dr_rms",    "resu_rms",   "resp_rms",   "resr_rms"
 1,  2.29801e-01,  5.03378e-02,  2.88455e-02,  1.07241e-02,  2.03448e-02,  2.03448e-02
 2,  2.58647e-01,  1.51207e-01,  1.17747e-01,  5.94100e-03,  5.40709e-03,  5.40709e-03
 3,  8.72627e-01,  5.95568e-01,  4.75178e-01,  4.18458e-03,  6.25604e-03,  6.25604e-03
 4,  3.39289e-01,  2.75104e-01,  2.11024e-01,  1.35157e-01,  2.57616e-01,  2.57616e-01
 5,  1.13501e+00,  5.81302e-01,  3.79980e-01,  1.81180e-01,  5.31819e-02,  5.31819e-02
 6,  4.67940e-01,  2.83047e-01,  2.98249e-01,  1.93246e-01,  3.24966e-01,  3.24966e-01
 7,  2.68325e-01,  1.12835e-01,  9.90863e-02,  5.66960e-02,  1.91540e-01,  1.91540e-01
 8,  2.00619e-01,  1.24138e-01,  7.90677e-02,  2.57842e-02,  6.30679e-02,  6.30679e-02
 9,  3.47459e-01,  2.28170e-01,  1.64274e-01,  8.91253e-03,  1.89085e-02,  1.89085e-02
\end{verbatim}

Note that the precision of the compressible file has been truncated for the purposes of 
illustrating the file format. The actual file has the same precision as the incompressible file.


\subsubsection{Flow field}
\label{subsubsec-form-flow}

As with the convergence histories, the format of the solution files depends on whether
incompressible or compressible flow has been modelled.
Sample files for both flow types are listed below, showing the solution with 8 axial
stations.

See \Cref{subsec-res-flow} for details on how to plot the flow field of a single run and
to compare the flow fields of multiple runs.


\begin{verbatim}
head nozzle_008_insi.sol 
"index", "x",        "area",    "density",  "velocity",  "pressure"
 1,  0.00000000,  1.00000000,  1.00000000,  0.99352315,  1.00645586
 2,  0.14285714,  0.95002040,  1.00000000,  1.05396039,  0.93273107
 3,  0.28571428,  0.91778253,  1.00000000,  1.09468920,  0.87557899
 4,  0.42857142,  0.90196380,  1.00000000,  1.12077187,  0.83702077
 5,  0.57142857,  0.90196380,  1.00000000,  1.12039273,  0.83109312
 6,  0.71428571,  0.91778253,  1.00000000,  1.10510955,  0.84489365
 7,  0.85714285,  0.95002040,  1.00000000,  1.06576098,  0.88937807
 8,  1.00000000,  1.00000000,  1.00000000,  1.01185327,  0.95000000


head nozzle_008_cosi.sol 
"index", "x",     "area", "density","velocity","pressure",    "temp",  "Mach no"
 1,  0.000000,  1.038230,  0.733619,  0.903193,  0.648126,  0.883463,  0.812124
 2,  0.142857,  1.002814,  0.673123,  1.019130,  0.569949,  0.851624,  0.933340
 3,  0.285714,  1.005860,  0.594401,  1.150608,  0.479892,  0.810871,  1.079907
 4,  0.428571,  1.039360,  0.524993,  1.260737,  0.405785,  0.772934,  1.211963
 5,  0.571428,  1.099747,  0.448136,  1.395860,  0.318806,  0.721653,  1.388716
 6,  0.714285,  1.185954,  0.441964,  1.312472,  0.369072,  0.753916,  1.277511
 7,  0.857142,  1.298492,  0.629288,  0.841891,  0.630211,  0.898745,  0.750540
 8,  1.000000,  1.438982,  0.791263,  0.604183,  0.750000,  0.947851,  0.524487

\end{verbatim}

As before, the precision of the output has been truncated for the purposes of illustrating 
the file format.


